% Source: https://artofproblemsolving.com/community/c3h2918242_argentina_nmo_round_2_intercollegiate_geo__level_1_200616
% Stable identity: aops-topic-2918242-q02
Let $\angle AOM=\alpha$. Then $\angle AON=\alpha+90^\circ$ and $\angle AOP=\dfrac{\alpha+90^\circ}2$. Since $OM$ bisects $\angle AOP$, $\alpha=\dfrac{\alpha+90^\circ}4$, so $\alpha=30^\circ$. Hence $\angle AON=120^\circ$ and $\angle BOM=150^\circ$, so $\angle BOQ=75^\circ$ and $\angle AOQ=105^\circ$.
\begin{asy}
usepackage("amsmath");
size(3.5cm);
defaultpen(fontsize(10pt));
draw(arc((0,0),1,0,180));draw((-1,0)--(1,0));pair O=(0,0);pair M=dir(150),N=dir(60),P=dir(120),Q=dir(75);
draw(O--M);draw(O--N);draw(O--P,dashed);draw(O--Q,dashed);label("$A$",(-1,0),W);label("$B$",(1,0),E);label("$M$",M,NW);label("$N$",N,NE);label("$P$",P,N);label("$Q$",Q,N);label("$O$",O,S);
\end{asy}

Therefore $\angle QON=120^\circ-105^\circ=\boxed{15^\circ}$.
